\documentclass[14pt,a4paper,oneside]{extbook}
\usepackage[utf8]{vietnam}
\usepackage{amsmath,amssymb,mathrsfs,fancyhdr,enumerate,multirow,makecell}
\usepackage{tkz-euclide,tikz-3dplot,tikz,tkz-tab}
\usetikzlibrary{arrows.meta,calc,intersections,angles,quotes,backgrounds,shapes.geometric,patterns,through,shadows,positioning}
\usepackage{yhmath}
\usepackage[top=0.7cm, bottom=0.7cm, left=1.0cm, right=1.0cm, headheight=16pt, headsep=3mm, footskip=5mm]{geometry}
\usepackage[hidelinks,unicode]{hyperref}
\usepackage{qrcode}
\usepackage[many]{tcolorbox}
\usepackage{fontawesome}

\pagestyle{fancy}
\fancyhf{}
\fancyhead[L]{\footnotesize\bfseries\color{teal!85!black}\faBook~SỔ TAY KIẾN THỨC TOÁN 4}
\fancyhead[R]{\footnotesize\bfseries\color{blue!85!black}Lớp Toán Thầy Lĩnh -- loptoanthaylinh.vn}
\fancyfoot[C]{\footnotesize\bfseries\color{gray!80!black}-- Sổ tay kiến thức: Bài 61. Phép trừ phân số --}
\renewcommand{\headrulewidth}{0.4pt}

\begin{document}

% --- BANNER TIÊU ĐỀ BÀI HỌC ---
\noindent
\begin{tcolorbox}[
  enhanced,
  colback=teal!8!white,
  colframe=teal!80!black,
  arc=3mm, boxrule=1pt,
  left=3mm, right=3mm, top=2mm, bottom=2mm,
  halign=center
]
  {\Large\bfseries\color{teal!90!black}\faBookmark\ BÀI 61: PHÉP TRỪ PHÂN SỐ}\\[1.5pt]
  {\footnotesize\itshape\color{black!75} Quy tắc trừ cùng mẫu, khác mẫu, tính chất và trừ số tự nhiên cho phân số trên 1 trang A4}
\end{tcolorbox}

\vspace{-2.5mm}

\noindent
\begin{minipage}[t]{0.49\textwidth}\vspace{0pt}
\begin{tcolorbox}[
  enhanced,
  colback=blue!4!white,
  colframe=blue!75!black,
  arc=2mm, boxrule=0.8pt,
  left=2mm, right=2mm, top=1.8mm, bottom=1.8mm,
  title={\footnotesize\bfseries\color{white}\faMinusCircle\ 1. Trừ hai phân số cùng mẫu}
]
{\footnotesize
\textbf{Quy tắc:} Muốn trừ hai phân số có cùng mẫu số, ta \textbf{trừ tử số của số bị trừ cho tử số của số trừ} và \textbf{giữ nguyên mẫu số}:
$$\dfrac{a}{m} - \dfrac{b}{m} = \dfrac{a - b}{m}$$
\textbf{Lưu ý:}
\begin{itemize}\setlength{\itemsep}{1pt}\setlength{\leftskip}{-2.5mm}
  \item Điều kiện trừ: Tử số bị trừ $\ge$ Tử số trừ ($a \ge b$).
  \item Mẫu số \textbf{giữ nguyên}, tuyệt đối không lấy mẫu trừ mẫu!
\end{itemize}
}
\end{tcolorbox}
\end{minipage}
\hfill
\begin{minipage}[t]{0.49\textwidth}\vspace{0pt}
\begin{tcolorbox}[
  enhanced,
  colback=teal!4!white,
  colframe=teal!75!black,
  arc=2mm, boxrule=0.8pt,
  left=2mm, right=2mm, top=1.8mm, bottom=1.8mm,
  title={\footnotesize\bfseries\color{white}\faRandom\ 2. Trừ hai phân số khác mẫu}
]
{\footnotesize
\textbf{Quy trình 2 bước chuẩn:}
\begin{itemize}\setlength{\itemsep}{1.5pt}\setlength{\leftskip}{-2.5mm}
  \item \textbf{Bước 1 (Quy đồng):} Quy đồng mẫu số hai phân số để có cùng mẫu số chung.
  \item \textbf{Bước 2 (Trừ tử):} Lấy tử số trừ tử số, giữ nguyên mẫu số chung.
  \item \textbf{Trừ số tự nhiên cho phân số:}\\
  $1 - \dfrac{a}{b} = \dfrac{b}{b} - \dfrac{a}{b} = \dfrac{b - a}{b}$; \quad $n - \dfrac{a}{b} = \dfrac{n \times b - a}{b}$.
  \item Luôn rút gọn kết quả về phân số tối giản.
\end{itemize}
}
\end{tcolorbox}
\end{minipage}

\vspace{-2.5mm}
% --- KHỐI 3: SƠ ĐỒ TƯ DUY TIKZ ---
\noindent
\begin{tcolorbox}[
  enhanced,
  colback=orange!4!white,
  colframe=orange!80!black,
  arc=2mm, boxrule=0.8pt,
  left=2mm, right=2mm, top=1.5mm, bottom=1.5mm,
  title={\footnotesize\bfseries\color{white}\faSitemap\ SƠ ĐỒ TƯ DUY: LƯU ĐỒ THỰC HIỆN PHÉP TRỪ PHÂN SỐ}
]
\centering
\begin{tikzpicture}[
  box/.style={draw, rounded corners=2pt, align=center, font=\scriptsize\bfseries, inner sep=3pt, minimum height=1.05cm},
  arr/.style={-{Stealth[length=2.5mm]}, thick}
]
  \node[box, fill=purple!12, draw=purple!80!black, text width=2.8cm] (root) at (0, 0) {PHÉP TRỪ\\$\dfrac{a}{b} - \dfrac{c}{d}$};

  \node[box, fill=blue!10, draw=blue!80!black, text width=4.0cm] (cung) at (4.8, 0.95) {CÙNG MẪU SỐ ($b = d$)\\Trừ tử, GIỮ NGUYÊN mẫu\\$\mathbf{\dfrac{a - c}{b}}$};

  \node[box, fill=teal!10, draw=teal!80!black, text width=4.0cm] (khac) at (4.8, -0.95) {KHÁC MẪU SỐ ($b \ne d$)\\Quy đồng mẫu số chung\\$\to$ Đưa về cùng mẫu};

  \node[box, fill=green!15, draw=green!70!black, text width=2.8cm] (res) at (10.6, 0) {RÚT GỌN KẾT QUẢ\\VỀ PHÂN SỐ\\TỐI GIẢN!};

  \draw[arr, blue!80!black] (root.east) -- ++(0.4, 0) |- (cung.west);
  \draw[arr, teal!80!black] (root.east) -- ++(0.4, 0) |- (khac.west);
  \draw[arr, green!70!black] (cung.east) -- ++(0.3, 0) |- (res.north);
  \draw[arr, green!70!black] (khac.east) -- ++(0.3, 0) |- (res.south);
\end{tikzpicture}
\end{tcolorbox}

\vspace{-2.5mm}
% --- KHỐI 4: VÍ DỤ & MẸO TRÁNH BẪY ---
\noindent
\begin{minipage}[t]{0.49\textwidth}\vspace{0pt}
\begin{tcolorbox}[
  enhanced,
  colback=green!4!white,
  colframe=green!60!black,
  arc=2mm, boxrule=0.8pt,
  left=2mm, right=2mm, top=1.8mm, bottom=1.8mm,
  title={\footnotesize\bfseries\color{white}\faCheckCircle\ Ví dụ thực hành chuẩn mực}
]
{\footnotesize
\textbf{1. Cùng mẫu:}\\
$\dfrac{5}{6} - \dfrac{1}{6} = \dfrac{5 - 1}{6} = \dfrac{4}{6} = \mathbf{\dfrac{2}{3}}$ (rút gọn về tối giản).\\
\textbf{2. Khác mẫu:}\\
$\dfrac{7}{5} - \dfrac{2}{3} = \dfrac{21}{15} - \dfrac{10}{15} = \mathbf{\dfrac{11}{15}}$ (MSC = 15).\\
$\dfrac{3}{4} - \dfrac{5}{12} = \dfrac{9}{12} - \dfrac{5}{12} = \dfrac{4}{12} = \mathbf{\dfrac{1}{3}}$ (MSC = 12).\\
\textbf{3. Trừ với số tự nhiên:}\\
$1 - \dfrac{3}{7} = \dfrac{7}{7} - \dfrac{3}{7} = \mathbf{\dfrac{4}{7}}; \quad 2 - \dfrac{3}{4} = \dfrac{8}{4} - \dfrac{3}{4} = \mathbf{\dfrac{5}{4}}$.
}
\end{tcolorbox}
\end{minipage}
\hfill
\begin{minipage}[t]{0.49\textwidth}\vspace{0pt}
\begin{tcolorbox}[
  enhanced,
  colback=red!4!white,
  colframe=red!70!black,
  arc=2mm, boxrule=0.8pt,
  left=2mm, right=2mm, top=1.8mm, bottom=1.8mm,
  title={\footnotesize\bfseries\color{white}\faExclamationTriangle\ Mẹo tránh bẫy trắc nghiệm}
]
{\footnotesize
\begin{itemize}\setlength{\itemsep}{1.5pt}\setlength{\leftskip}{-2.5mm}
  \item \textbf{Bẫy lấy mẫu trừ mẫu:}\\
  Ví dụ $\frac{5}{6} - \frac{1}{6}$ học sinh hay viết nhầm mẫu thành $6 - 6 = 0$ (SAI nghiêm trọng!). Mẫu số PHẢI GIỮ NGUYÊN.
  \item \textbf{Bẫy đổi số 1 thành phân số:}\\
  Khi gặp $1 - \frac{a}{b}$, nhanh trí đổi $1 = \frac{b}{b}$ (bằng chính mẫu số của phân số trừ) để trừ ngay lập tức.
\end{itemize}
}
\end{tcolorbox}
\end{minipage}

\end{document}
